\chapter{Disponibilidade dos Dados e Recomposição da Evidência}
\label{apend:recuperacao}

Parte das evidências de apoio desta dissertação repousa em artefatos processados cujos dados brutos não acompanham o repositório. Este apêndice registra quais são, por que não estão versionados, e o que é --- e o que não é --- possível refazer a partir do que está distribuído.

\section{O que não está versionado, e por quê}

Três categorias de dados ficam fora do controle de versão, por razões distintas:

\begin{itemize}
\item \textbf{Saídas brutas da plataforma de otimização.} Cada execução produz um arquivo binário com a população final e o histórico de avaliações. Para as 51 instâncias sintéticas, com nove algoritmos e 60 execuções por configuração, esse conjunto excede em ordens de grandeza os limites práticos de um repositório de código.

\item \textbf{Trajetórias por geração.} As execuções instrumentadas que registram observáveis a cada geração --- usadas na análise de dinâmica populacional --- somam alguns gigabytes.

\item \textbf{Dados intermediários das fases de calibração e ablação.} As varreduras de configuração produziram milhares de execuções cujas saídas foram consolidadas em artefatos resumidos; os intermediários não foram preservados no repositório.
\end{itemize}

O que \emph{está} versionado, em todos os casos, é o produto processado: as tabelas de métricas por execução, os resumos estatísticos e os artefatos de auditoria a partir dos quais as tabelas desta dissertação são geradas. As contagens reportadas foram conferidas contra esses artefatos.

\section{Três estados distintos, e por que a distinção importa}

Trabalhos que declaram disponibilidade de dados costumam tratar repositório e depósito como sinônimos. Aqui eles não são, e confundi-los produziria uma promessa que esta pesquisa não pode cumprir. Três estados precisam ser separados:

\begin{itemize}
\item \textbf{Repositório de código.} Mantém a implementação do algoritmo, os roteiros de experimento, os scripts de análise, os dados processados e os geradores das tabelas desta dissertação. É o estado vivo, e é dele que a Seção~\ref{sec:disponibilidade} dá o endereço.

\item \textbf{Registro de arquivamento.} Preserva instantâneos versionados desse repositório, cada um com identificador próprio e permanente. As versões publicadas até a redação desta dissertação correspondem aos manuscritos derivados da pesquisa; o identificador de conceito resolve sempre para a mais recente.

\item \textbf{Fora de ambos.} As saídas brutas de execução. Nenhuma versão publicada do registro as contém, e o repositório as exclui deliberadamente pelos motivos da seção anterior.
\end{itemize}

A consequência é direta e é registrada aqui em vez de ser deixada para o leitor descobrir: \textbf{não há depósito de onde recuperar os dados brutos}. O que está arquivado é a cadeia processada, não a sua origem.

\section{Consequência para a reprodutibilidade}

A distinção entre reproduzir, regenerar e reexecutar é relevante aqui, e esta dissertação não a dissimula.

As análises da família confirmatória são \textbf{regeneráveis} a partir dos dados processados versionados: os scripts declarados no Apêndice~\ref{apend:fontes} reconstroem as contagens e as tabelas a partir do repositório, sem depender de dados externos.

As famílias de calibração, ablação, transferência para engenharia e do controlador são \textbf{verificáveis, mas não regeneráveis} no estado atual do repositório: os seus artefatos processados estão presentes e foram conferidos, mas os produtores não executam a partir de um clone limpo, porque os seus dados de entrada não estão versionados.

Para essas famílias, a via disponível não é recuperar, é \textbf{reexecutar}. Os roteiros de experimento que produziram as execuções originais estão versionados, e reconstruir os dados de entrada exige rodá-los novamente sob a mesma plataforma, com os mesmos identificadores de execução e a mesma configuração declarada no Capítulo~\ref{sec:Experiments}. O custo é o da campanha experimental original, e não o de um download. O repositório distribui o que é necessário para refazer o percurso; não distribui um atalho para o seu resultado intermediário.

Essa assimetria está declarada nas ameaças à validade, no Capítulo~\ref{sec:threats}, e é registrada aqui para que o leitor possa distinguir, ao consultar cada resultado, qual dos três níveis de garantia se aplica.

\section{Declaração de disponibilidade de dados e código}
\label{sec:disponibilidade}

O código do algoritmo, os roteiros de experimento, os scripts de análise, os dados processados e os geradores das tabelas desta dissertação são mantidos no repositório público \url{https://github.com/pedsanches/ivfspea2-content} e arquivados no Zenodo sob o DOI de conceito \texttt{10.5281/zenodo.19071253}~\cite{ivfspea2github}. O DOI de conceito resolve sempre para a versão mais recente do registro, e é por isso que ele --- e não o identificador de uma versão específica --- é o citado aqui.

A implementação avaliada é a classe \texttt{IVFSPEA2V2}, no diretório de algoritmos multiobjetivo da distribuição do PlatEMO incluída no repositório. O identificador é nome de classe na implementação, e não designação de variante do método descrito no Capítulo~\ref{sec:Proposal}.

O registro arquiva, além do código, os manifestos que associam cada conjunto de dados processados ao script que o produziu e à coorte de execuções que o originou. É esse encadeamento --- e não a presença dos dados brutos, que o registro não contém --- que permite verificar se um número reportado corresponde ao artefato que o sustenta. O Apêndice~\ref{apend:fontes} relaciona, fonte a fonte, qual artefato sustenta cada afirmação do texto.
